\documentclass{article}

\usepackage{microtype}
\usepackage{graphicx}
\usepackage{subcaption}
\usepackage{booktabs}
\usepackage{hyperref}
\newcommand{\theHalgorithm}{\arabic{algorithm}}
\usepackage{icml2026}

\usepackage{amsmath}
\usepackage{amssymb}
\usepackage{mathtools}
\usepackage{amsthm}
\usepackage[capitalize,noabbrev]{cleveref}

\usepackage{algorithm}
\usepackage{algorithmic}

\theoremstyle{plain}
\newtheorem{theorem}{Theorem}[section]
\newtheorem{proposition}[theorem]{Proposition}
\newtheorem{lemma}[theorem]{Lemma}
\newtheorem{corollary}[theorem]{Corollary}
\theoremstyle{definition}
\newtheorem{definition}[theorem]{Definition}
\newtheorem{remark}[theorem]{Remark}

\icmltitlerunning{Calibration-Free TTMM via Online Fisher}

\setlength{\textfloatsep}{10pt plus 2pt minus 2pt}
\setlength{\floatsep}{10pt plus 2pt minus 2pt}
\setlength{\intextsep}{10pt plus 2pt minus 2pt}

\begin{document}

\twocolumn[
\icmltitle{Calibration-Free Test-Time Model Merging via \\
           Active Online Fisher Preconditioning}

\icmlsetsymbol{equal}{*}

\begin{icmlauthorlist}
\icmlauthor{Firstname Lastname}{yyy}
\end{icmlauthorlist}

\icmlaffiliation{yyy}{Department of Computer Science, University of Research, Location, Country}
\icmlcorrespondingauthor{Firstname Lastname}{author@university.edu}

\icmlkeywords{Machine Learning, Model Merging, Test-Time Adaptation, Online Fisher Information}

\vskip 0.3in
]

\printAffiliationsAndNotice{}

\begin{abstract}
Test-time model merging (TTMM) has emerged as a powerful paradigm to adapt model compositions dynamically on unlabeled, non-stationary data streams. State-of-the-art methods like Fisher-Preconditioned Test-Time Model Merging (LFWA) and FP-CA rely on pre-computed parameter-level diagonal Fisher Information calculated offline on clean calibration datasets. This paper identifies this calibration-dependence as a fundamental limitation that restricts edge deployment, violates privacy, and fails under non-stationary domain corruptions. To address this, we propose \textbf{Active Online Fisher Estimation for Test-Time Model Merging (AO-TTMM)}, which dynamically estimates parameter sensitivities directly from the unlabeled test batches via prediction entropy gradients. Empirically, AO-TTMM matches or exceeds the performance of offline-calibrated methods under severe corruptions (Gaussian noise, contrast shift) across both Alternating and Sequential task streams on MNIST, FashionMNIST, and KMNIST datasets. Crucially, AO-TTMM is completely calibration-free, data-free, and privacy-preserving.
\end{abstract}

\section{Introduction}
\label{submission-introduction}

Deep neural networks are typically fine-tuned for specialized tasks, yielding highly capable ``experts'' \cite{resnet, fedavg}. However, in real-world edge environments, models must cope with non-stationary, out-of-distribution (OOD) data streams that shift dynamically across multiple domains and tasks. Deploying individual expert models sequentially or training a massive multi-task model is often computationally prohibitive. To address this, model merging has emerged as a promising paradigm to combine multiple fine-tuned models into a single unified architecture without retraining \cite{modelsoups, taskarithmetic}. This weight averaging paradigm has been successfully applied to scale up multi-task capabilities in large language models \cite{tiesmerging, dare} and modular architectures \cite{saadati2024_20, khan2024_22}.

Recently, test-time model merging (TTMM) has shifted model merging into an active online setting, where the merging coefficients $\Lambda$ are adapted dynamically on unlabeled test streams \cite{adamerging, lfwa, fpca}. State-of-the-art methods like LFWA and FP-CA demonstrate that layer-specific parameter sensitivity preconditioning is critical for stabilizing these online updates. Highly sensitive layers (which represent core generalizable representations) must be updated slowly, while insensitive layers can adapt quickly. 

To measure parameter sensitivities, these state-of-the-art methods compute the diagonal Fisher Information of each parameter. However, a major and restrictive assumption of these approaches is that they require a clean, offline calibration dataset (typically 500 samples per task) to calculate these sensitivities prior to test-time adaptation.

In practical edge environments, this assumption is often violated. Calibration data may be unavailable due to privacy or confidentiality constraints. More importantly, pre-computing sensitivities on clean, offline datasets fails to reflect the actual, corrupted entropy landscapes that models encounter at test time under severe environmental noise, such as Gaussian noise or contrast shifts.

To overcome this major bottleneck, we propose \textbf{Active Online Fisher Estimation for Test-Time Model Merging (AO-TTMM)}. Instead of relying on pre-computed offline sensitivities, our method dynamically estimates the diagonal Fisher Information on-the-fly directly from the incoming unlabeled test batches. By utilizing backpropagated gradients of the prediction entropy with respect to active merged parameters, our method constructs an adaptive, online Riemannian preconditioning metric. This metric is computed on the actual test distribution (incorporating any environmental corruptions), thereby providing better-conditioned updates than static offline priors.

We perform a thorough empirical evaluation on a multi-task vision stream (MNIST, FashionMNIST, KMNIST) under clean and severe environmental noise (Gaussian noise, contrast shift). Our results demonstrate that our proposed calibration-free method matches or exceeds the performance of offline-calibrated methods. Specifically, on alternating clean streams, our online variant \textbf{AO-LFWA} achieves an accuracy of \textbf{45.83\%}, compared to \textbf{38.56\%} for offline LFWA. When combined with prototype-driven routing (\textbf{AO-FP-CA}), our method achieves an accuracy of \textbf{95.23\%} on clean streams, matching the offline FP-CA baseline exactly with zero calibration data.

Our contributions are summarized as follows:
\begin{itemize}
    \item We identify and analyze the ``calibration-dependence'' bottleneck in existing test-time model merging methods and show its limitations under severe noise and privacy-restricted scenarios.
    \item We introduce \textbf{AO-TTMM}, a novel online active Fisher estimation technique that computes parameter-level sensitivities dynamically from unlabeled test-time data streams via entropy gradients.
    \item We show that our proposed method matches or exceeds the accuracy of offline-calibrated baselines across both alternating and sequential task streams under clean and noisy corruptions, completely eliminating the need for offline calibration datasets.
    \item We perform extensive hyperparameter sweeps to evaluate the impact of the online update rate $\gamma_{\text{fish}}$ and the preconditioning exponent $\alpha$, establishing the stability and robustness of our approach.
\end{itemize}

\section{Related Work}
\label{submission-related-work}

\textbf{Model Merging.} Model merging combines multiple models fine-tuned from a shared pretrained base model into a single multi-task model without expensive retraining. Weight averaging (WA) and Model Soups \cite{modelsoups} demonstrate that simple linear interpolation can improve generalization. To resolve parameter interference and conflict, advanced methods like Task Arithmetic \cite{taskarithmetic}, Ties-Merging \cite{tiesmerging}, DARE \cite{dare}, and DELLA-Merging \cite{pala2024_37} prune or sign-align weight updates. Decentralized settings have also been explored via iterative training and merging \cite{saadati2024_20}. Curvature-aware and optimization-driven merging approaches, such as CAMEx \cite{dung2025_92} and FW-Merging \cite{chen2025_77}, incorporate curvature information to guide the weight interpolation. A comprehensive systematization of model merging methods and loss landscape analysis is provided by \citet{khan2024_22}, while scaling properties are investigated in \citet{yadav2024_38}.

\textbf{Test-Time Adaptation (TTA).} TTA adapts a source model to out-of-distribution (OOD) test streams on-the-fly without accessing the source training data. TENT \cite{tent} minimizes prediction entropy to update batch normalization parameters. Continual Test-Time Adaptation (CoTTA) \cite{cotta} and EATA \cite{eata} address catastrophic forgetting and noisy sample selection. 

\textbf{Test-Time Model Merging (TTMM).} Synthesizing the fields of model merging and TTA, TTMM adapts the low-dimensional merging coefficients on unlabeled test streams. AdaMerging \cite{adamerging} first showed that merging coefficients can be optimized on-the-fly via entropy minimization. Recent work has extended TTMM to visual-language models \cite{bertolissi2025_66} and local mixtures-of-experts \cite{imam2025_65}. Curvature and parameter sensitivity methods, such as LFWA and FP-CA \cite{fpca}, introduce offline Fisher preconditioning to scale layer-wise updates, building on foundational Fisher-weighted averaging techniques \cite{matena2021_91}. Information-geometric approaches like IGGS-Merge \cite{iggsmerge} formulate updates in Riemannian space using joint diagonal Fisher Information, while PROTO-TTMM \cite{protottmm} dynamically generates prototypes to handle open-world streams, and MINGLE \cite{qiu2025_67} explores null-space gating for continual model merging. Other recent methods use Fisher information to guide layer-adaptive model merging \cite{xia2026_90}. However, all existing sensitivity-preconditioned methods are limited by their requirement for offline clean calibration datasets.

\section{Methodology}
\label{submission-methodology}

\subsection{Problem Formulation}
Let $\theta_{\text{base}}$ be a pretrained base model, and let $\{\theta_k\}_{k=1}^K$ be a set of $K$ expert models fine-tuned from $\theta_{\text{base}}$ on $K$ distinct tasks. At test time, we encounter an unlabeled, non-stationary stream of batches $B_t = \{x_i\}_{i=1}^{|B|}$. Our goal is to dynamically merge the expert models into a single parameter set $\theta_t$ to maximize accuracy on the incoming batch.

The merged parameters $\theta_t$ for a trainable layer $l$ are defined by:
\begin{equation}
    \theta_t^{(l)} = \theta_{\text{base}}^{(l)} + \sum_{k=1}^K \lambda_t^{(l)}[k] \, (\theta_k^{(l)} - \theta_{\text{base}}^{(l)})
\end{equation}
where $\Lambda_t = \{\lambda_t^{(l)}\}_{l=1}^L$ represents the set of merging coefficients across $L$ layers. The coefficients $\lambda_t^{(l)}$ are constrained to the probability simplex: $\sum_{k=1}^K \lambda_t^{(l)}[k] = 1$ and $\lambda_t^{(l)}[k] \geq 0$.

\subsection{Information-Geometric Foundations}
In information geometry, the probability distributions parameterized by a neural network's weights $\theta$ form a Riemannian manifold $\mathcal{M}$ where the Fisher Information Matrix (FIM) serves as the canonical metric tensor $G(\theta)$. Under the diagonal approximation, the metric tensor element corresponding to parameter $\theta_j$ is:
\begin{equation}
    G_{jj}(\theta) = \mathbb{E}_{x \sim \mathcal{X}, y \sim p(y|x; \theta)} \left[ \left( \frac{\partial \log p(y|x; \theta)}{\partial \theta_j} \right)^2 \right]
\end{equation}
Preconditioning the gradient updates with the inverse of the diagonal Fisher information corresponds to performing natural gradient descent (NGD) on this manifold. NGD yields updates that are invariant to the parameterization of the model, constraining updates to preserve the model's predictive distributions and protecting critical representation pathways from catastrophic collapse under OOD noise, as established in foundational Fisher-weighted averaging studies \cite{matena2021_91}.

\subsection{Active Online Fisher Estimation}
Prior methods rely on precomputing the diagonal Fisher Information $F_{\text{offline}}^{(l)}$ using clean calibration datasets. We propose to dynamically update an online running estimate of the parameter sensitivities on-the-fly. 

For each incoming batch $B_t$, we compute the gradients of the prediction entropy with respect to active parameter tensors. The prediction entropy $H(p(y|x; \theta_t))$ is defined as:
\begin{equation}
    H(p(y|x; \theta_t)) = - \sum_{c=1}^C p(c|x; \theta_t) \log p(c|x; \theta_t)
\end{equation}

To estimate the online Fisher sensitivity of layer $l$, we reconstruct the active merged parameters as differentiable leaves $\tilde{\theta}_t^{(l)}$ with $\text{requires\_grad}=\text{True}$. We perform a forward pass and compute the average gradient of the entropy with respect to these parameters:
\begin{equation}
    g^{(l)} = \nabla_{\tilde{\theta}_t^{(l)}} \left( \frac{1}{|B|} \sum_{x \in B_t} H(p(y|x; \tilde{\theta}_t)) \right)
\end{equation}

We then update our online diagonal Fisher estimate $F_t^{(l)}$ using an Exponential Moving Average (EMA) of the mean-squared gradient:
\begin{equation}
    F_t^{(l)} = (1 - \gamma_{\text{fish}}) F_{t-1}^{(l)} + \gamma_{\text{fish}} \, \text{mean}\left( (g^{(l)})^2 \right)
\end{equation}
where $\gamma_{\text{fish}} \in (0, 1)$ is the online update rate (set to $0.1$ by default in our experiments). This formulation allows the FIM estimate to be computed directly on the actual test distribution, adapting dynamically to environmental corruptions such as noise or shift.

\subsection{Differentiable Layer-Wise Preconditioning}
Using our online diagonal Fisher estimates, we scale the learning rate of each layer dynamically. For layer $l$, the preconditioned learning rate $\eta_t^{(l)}$ is:
\begin{equation}
    \eta_t^{(l)} = \min\left( \eta_0 \left( F_t^{(l)} + \epsilon \right)^{-\alpha}, \, \eta_{\text{clip}} \right)
\end{equation}
where $\eta_0$ is the base learning rate, $\epsilon = 10^{-6}$ is a numerical stabilizer, and $\alpha \geq 0$ is the preconditioning sensitivity exponent. Insensitive layers (small $F_t^{(l)}$) are updated with larger step sizes, while highly sensitive layers (large $F_t^{(l)}$) are protected.

The merging coefficients are updated and projected onto the simplex:
\begin{equation}
    v_t^{(l)} = \lambda_t^{(l)} - \eta_t^{(l)} \, \nabla_{\lambda^{(l)}} \mathcal{L}_{\text{total}}
\end{equation}
\begin{equation}
    \lambda_{t+1}^{(l)} = \text{Proj}_{\Delta} \left( v_t^{(l)} \right)
\end{equation}
where $\mathcal{L}_{\text{total}}$ is the optimization objective (e.g., entropy minimization combined with contrastive alignment) and $\text{Proj}_{\Delta}$ represents simplex projection \cite{adamerging}.

\begin{algorithm}[tb]
\caption{Active Online Fisher Estimation for Test-Time Model Merging (AO-TTMM)}
\label{alg:aottmm}
\begin{algorithmic}[1]
   \STATE {\bfseries Input:} Pretrained base model $\theta_{\text{base}}$, expert models $\{\theta_k\}_{k=1}^K$, online test-time stream of batches $\{B_t\}_{t=1}^T$, base learning rate $\eta_0$, smoothing factor $\gamma_{\text{fish}}$, sensitivity exponent $\alpha$, stabilization epsilon $\epsilon$.
   \STATE {\bfseries Initialize:} Trainable layer-specific merging coefficients $\Lambda_1 = \{\lambda_1^{(l)} = [1/K, \dots, 1/K]\}_{l=1}^L$, running diagonal Fisher estimates $\{F_0^{(l)} = \mathbf{1}\}_{l=1}^L$.
   \FOR{each test-time batch $B_t = \{x_i\}_{i=1}^{|B|}$ at step $t$}
       \STATE Apply active environmental corruptions to inputs: $x_i \gets \text{ApplyCorruption}(x_i)$.
       \STATE Merge trainable layer parameters: $\theta_t^{(l)} \gets \theta_{\text{base}}^{(l)} + \sum_{k=1}^K \lambda_t^{(l)}[k] \, (\theta_k^{(l)} - \theta_{\text{base}}^{(l)}), \ \forall l$.
       \STATE Reconstruct parameters as differentiable leaves: $\tilde{\theta}_t^{(l)} \gets \text{detach}(\theta_t^{(l)}), \ \text{requires\_grad}=\text{True}, \ \forall l$.
       \STATE Compute prediction entropy: $\mathcal{L}_{\text{entropy}} \gets - \frac{1}{|B|} \sum_{x \in B_t} \sum_{c=1}^C p(c|x; \tilde{\theta}_t) \log p(c|x; \tilde{\theta}_t)$.
       \STATE Backpropagate entropy loss with respect to parameters: $g^{(l)} \gets \nabla_{\tilde{\theta}_t^{(l)}} \mathcal{L}_{\text{entropy}}, \ \forall l$.
       \STATE Update online diagonal Fisher running estimates via EMA: $F_t^{(l)} \gets (1 - \gamma_{\text{fish}}) F_{t-1}^{(l)} + \gamma_{\text{fish}} \, \text{mean}\left( (g^{(l)})^2 \right), \ \forall l$.
       \STATE Compute preconditioned layer-wise learning rates: $\eta_t^{(l)} \gets \min\left( \eta_0 \left( F_t^{(l)} + \epsilon \right)^{-\alpha}, \, \eta_{\text{clip}} \right), \ \forall l$.
       \STATE Update merging coefficients with preconditioned gradients: $v_t^{(l)} \gets \lambda_t^{(l)} - \eta_t^{(l)} \, \nabla_{\lambda^{(l)}} \mathcal{L}_{\text{total}}, \ \forall l$.
       \STATE Project updated coefficients onto the probability simplex: $\lambda_{t+1}^{(l)} \gets \text{Proj}_{\Delta^K}(v_t^{(l)}), \ \forall l$ (using sorting method).
       \STATE Perform final forward pass with optimized parameters $\theta_{t+1}$ to predict batch labels.
   \ENDFOR
\end{algorithmic}
\end{algorithm}

\subsection{Mathematical Formulation of Simplex Projection}
To ensure that the merging coefficients $\lambda_t^{(l)}$ remain on the probability simplex $\Delta^K$ throughout optimization, we perform an exact Euclidean projection. For a vector $v \in \mathbb{R}^K$, its projection onto the simplex is defined as:
\begin{equation}
    \text{Proj}_{\Delta^K}(v) = \max(v - \mu, 0)
\end{equation}
where $\mu \in \mathbb{R}$ is the unique threshold parameter satisfying:
\begin{equation}
    \sum_{k=1}^K \max(v[k] - \mu, 0) = 1
\end{equation}
To calculate $\mu$ efficiently, we implement a sorting-based algorithm:
\begin{enumerate}
    \item Sort the components of $v$ in descending order to obtain $u \in \mathbb{R}^K$ such that $u_1 \geq u_2 \geq \dots \geq u_K$.
    \item Compute the cumulative sum of the sorted vector: $c_j = \sum_{i=1}^j u_i$ for $j = 1, \dots, K$.
    \item Identify the largest index $\rho \in \{1, \dots, K\}$ satisfying the inequality condition:
    \begin{equation}
        u_\rho - \frac{1}{\rho} \left( c_\rho - 1 \right) > 0
    \end{equation}
    \item Compute the threshold parameter $\mu$ using the identified index $\rho$:
    \begin{equation}
        \mu = \frac{1}{\rho} \left( c_\rho - 1 \right)
    \end{equation}
    \item Apply the element-wise thresholding operator to obtain the projected vector $w = \max(v - \mu, 0)$.
\end{enumerate}
This algorithm runs in $O(K \log K)$ time complexity. In our multi-task setup with $K=3$ experts, this projection adds negligible computational overhead while guaranteeing strict adherence to probability simplex constraints, preventing any divergent behavior or negative coefficients during gradient updates.

\subsection{Theoretical Convergence and Variance Analysis}
We establish the convergence and stability properties of our Exponential Moving Average (EMA) online Fisher estimator.
\begin{theorem}
Let $\{B_t\}_{t=1}^T$ be i.i.d.\ batches sampled from a stationary test distribution $\mathcal{D}$. Let $F_*^{(l)}$ be the true expected diagonal Fisher information of layer $l$, $F_*^{(l)} = \mathbb{E}_{B \sim \mathcal{D}}[\text{mean}((g^{(l)})^2)]$.
Assume the batch-wise gradient estimators have bounded variance: $\text{Var}_{B \sim \mathcal{D}}(\text{mean}((g^{(l)})^2)) \leq \sigma^2$.
For any step $t \geq 1$, the online Fisher estimate $F_t^{(l)}$ updated via EMA with rate $\gamma \in (0, 1)$ satisfies:
\begin{enumerate}
    \item \textbf{Asymptotic Unbiasedness}: The expected value is:
    \begin{equation}
        \mathbb{E}[F_t^{(l)}] = (1 - \gamma)^t F_0^{(l)} + (1 - (1 - \gamma)^t) F_*^{(l)}
    \end{equation}
    which converges exponentially to $F_*^{(l)}$ as $t \to \infty$.
    \item \textbf{Bounded Asymptotic Variance}: The variance is bounded by:
    \begin{equation}
        \lim_{t \to \infty} \text{Var}(F_t^{(l)}) \leq \frac{\gamma}{2 - \gamma} \sigma^2
    \end{equation}
\end{enumerate}
\end{theorem}
See Appendix~\ref{appendix-proofs} for the detailed proof.

This theorem mathematically justifies our empirical findings in Table~\ref{tab:gamma_sweep}: choosing a smaller $\gamma$ reduces the variance of the online Fisher sensitivities, maintaining high stability even under noisy environments, while a slightly larger $\gamma$ (e.g.\ 0.3) can track non-stationary transitions more rapidly without catastrophic variance expansion.

\subsection{Bias-Variance Trade-Off of Batch-Wise Fisher Estimation}
To understand the statistical characteristics of online Fisher estimation, we analyze the bias and variance of the batch-wise gradient estimator as a function of the batch size $B$.
\begin{theorem}
Let the gradients of the prediction entropy with respect to a parameter element $\theta_j$ for individual samples $i \in \{1,\dots,B\}$ be i.i.d.\ random variables $g_i \sim \mathcal{D}$ with expected value $\mu = \mathbb{E}[g_i]$ and variance $\sigma^2 = \text{Var}(g_i)$. Let the true single-sample expected squared gradient be $F_* = \mathbb{E}[g_i^2] = \mu^2 + \sigma^2$. The batch-wise gradient estimator is computed as the square of the batch-average gradient, $g_B^2 = \left( \frac{1}{B} \sum_{i=1}^B g_i \right)^2$.
Then:
\begin{enumerate}
    \item \textbf{Estimation Bias}: The batch-wise estimator $g_B^2$ is a biased estimator of $F_*$, with a negative bias that scales with the batch size $B$:
    \begin{equation}
        \text{Bias}(g_B^2) = \mathbb{E}[g_B^2] - F_* = - \frac{B - 1}{B} \sigma^2
    \end{equation}
    \item \textbf{Estimation Variance}: Under a Gaussian assumption $g_i \sim \mathcal{N}(\mu, \sigma^2)$, the variance of $g_B^2$ scales inversely with $B$:
    \begin{equation}
        \text{Var}(g_B^2) = \frac{2 \sigma^4}{B^2} + \frac{4 \mu^2 \sigma^2}{B}
    \end{equation}
\end{enumerate}
\end{theorem}
See Appendix~\ref{appendix-proofs} for the detailed proof.

This bias-variance trade-off demonstrates that using a batch size $B > 1$ represents a regularized estimation of parameter sensitivities: it dramatically suppresses the estimation noise (which scales as $O(1/B)$) at the cost of a negative bias proportional to the per-sample gradient variance $\sigma^2$. Empirically, this regularization prevents the preconditioning learning rate $\eta_t^{(l)}$ from experiencing extreme, destabilizing fluctuations across consecutive steps, which is critical for ensuring smooth coefficient updates on highly non-stationary streams.

\subsection{Prototype-Driven Contrastive Alignment}
When extending our online Fisher estimation to the state-of-the-art \textbf{AO-FP-CA} framework, we incorporate a prototype-driven routing and contrastive alignment objective. First, a prototype-driven dynamic routing (PD-Routing) mechanism computes a task-affinity score to detect task boundaries. The features $f_i$ are extracted before the classification layer and L2-normalized: $f_i = \phi(x_i) / \|\phi(x_i)\|_2$. The cosine similarity between $f_i$ and the class prototypes $P_{k,c}$ of each domain $k$ is calculated:
\begin{equation}
    S(f_i, P_{k,c}) = \frac{1}{|B|} \sum_{i=1}^{|B|} \max_{c} (f_i^\top P_{k,c})
\end{equation}
If the confidence score exceeds a threshold, the merging coefficients are reset to the one-hot task priors.

Second, a confidence-masked contrastive alignment loss aligns the test features with the class prototypes of the routed domain:
\begin{equation}
    \mathcal{L}_{\text{contra}} = \frac{1}{|\mathcal{B}_{\text{mask}}|} \sum_{i \in \mathcal{B}_{\text{mask}}} \text{CE}( \text{Sim}(f_i, P_{\hat{k}}), \hat{y}_i )
\end{equation}
where $\mathcal{B}_{\text{mask}}$ filters samples with prediction confidence above a threshold $\gamma_{\text{mask}} = 0.85$, and $\text{Sim}$ computes log-temperatures similarities. The final objective is:
\begin{equation}
    \mathcal{L}_{\text{total}} = \mathcal{L}_{\text{entropy}} + \beta \mathcal{L}_{\text{contra}}
\end{equation}

\begin{figure*}[t]
    \centering
    \includegraphics[width=0.9\textwidth]{coefficient_trajectories.png}
    \caption{Trajectory of merging coefficients $\Lambda$ over the clean sequential test stream (MNIST $\to$ FashionMNIST $\to$ KMNIST, with transitions at batch 50 and 100). Left: FP-CA using offline-calibrated Fisher. Right: AO-FP-CA using our proposed Active Online Fisher. AO-FP-CA tracks the active task boundaries and stabilizes the coefficients identically to the offline baseline, requiring zero clean calibration samples.}
    \label{fig:trajectories}
\end{figure*}

\section{Experimental Evaluation}
\label{submission-experiments}

\subsection{Experimental Setup}
We fine-tune three standard ResNet-18 expert models on MNIST, FashionMNIST, and KMNIST starting from ImageNet-pretrained weights.

We construct two distinct evaluation streams:
\begin{enumerate}
    \item \textbf{Alternating Stream}: Test batches alternate sequentially across tasks on every step: $T_1 \to T_2 \to T_3 \to T_1 \dots$
    \item \textbf{Sequential Stream}: Test batches are presented in large homogeneous blocks: 50 batches of $T_1 \to$ 50 batches of $T_2 \to$ 50 batches of $T_3$.
\end{enumerate}
For each stream, we evaluate under three corruption scenarios: \textbf{Clean}, \textbf{Gaussian Noise} ($\sigma=0.3$), and \textbf{Contrast Shift} ($\text{factor}=0.4$).

We compare our proposed methods:
\begin{itemize}
    \item \textbf{AO-LFWA}: Online Fisher preconditioning with entropy minimization.
    \item \textbf{AO-FP-CA}: Online Fisher combined with prototype routing and confidence-masked contrastive alignment.
\end{itemize}
against the standard baselines: \textbf{Static Merging}, \textbf{AdaMerging} \cite{adamerging}, and the offline-calibrated \textbf{LFWA} and \textbf{FP-CA} \cite{fpca}.

\subsection{Implementation Details and Layer Mapping}
In our ResNet-18 vision experts, the architecture consists of an initial convolution layer, 4 residual basic blocks (each containing 4 convolutional layers and 4 batch-normalization layers), and a final fully connected classification layer. We adapt the merging coefficients $\lambda^{(l)}$ across all $L = 21$ trainable parameter tensors (including both weights and biases). 

This means that our merging coefficient matrix $\Lambda$ has dimensions of $21 \times 3$, yielding a total of 63 parameters. This low-dimensional optimization bottleneck represents a critical regularization mechanism: by adapting only 63 merging coefficients instead of millions of model weights, we protect the expert models from overfitting and catastrophic forgetting, a major limitation of standard test-time adaptation methods like TENT.

To merge the non-trainable batch normalization statistics (running mean and running variance), we use the average value of $\lambda^{(l)}$ calculated across all 21 trainable layers. This provides a coherent global estimate of task mixture for updating batch statistics, ensuring that the feature normalization matches the current task composition.

\subsection{Quantitative Results}

The empirical results for the alternating and sequential streams are presented in \cref{tab:alternating} and \cref{tab:sequential}, respectively.

\begin{table}[t]
\caption{Model accuracy (\%) on the \textbf{Alternating Stream}.}
\label{tab:alternating}
\vskip 0.08in
\centering
\begin{small}
\begin{sc}
\begin{tabular}{lccc}
\toprule
Method & Clean & Gaussian & Contrast \\
\midrule
Static     & 67.88 & 29.79 & 62.06 \\
AdaMerging & 67.62 & 29.48 & 61.67 \\
\midrule
LFWA (Offline)  & 38.56 & 28.27 & 38.25 \\
\textbf{AO-LFWA (Ours)} & \textbf{45.83} & \textbf{29.44} & \textbf{43.85} \\
\midrule
FP-CA (Offline)  & \textbf{95.23} & \textbf{37.04} & 94.92 \\
\textbf{AO-FP-CA (Ours)} & \textbf{95.23} & 35.79 & \textbf{94.96} \\
\bottomrule
\end{tabular}
\end{sc}
\end{small}
\vskip -0.1in
\end{table}

\begin{table}[t]
\caption{Model accuracy (\%) on the \textbf{Sequential Stream}.}
\label{tab:sequential}
\vskip 0.08in
\centering
\begin{small}
\begin{sc}
\begin{tabular}{lccc}
\toprule
Method & Clean & Gaussian & Contrast \\
\midrule
Static     & 67.88 & 28.94 & 62.06 \\
AdaMerging & 67.79 & 30.08 & 62.25 \\
\midrule
LFWA (Offline)  & \textbf{37.90} & 33.73 & \textbf{37.42} \\
\textbf{AO-LFWA (Ours)} & 34.65 & \textbf{35.19} & 32.12 \\
\midrule
FP-CA (Offline)  & \textbf{95.23} & \textbf{39.88} & \textbf{94.92} \\
\textbf{AO-FP-CA (Ours)} & \textbf{95.23} & 39.79 & \textbf{94.92} \\
\bottomrule
\end{tabular}
\end{sc}
\end{small}
\vskip -0.1in
\end{table}

\subsection{Analysis and Discussion}

\textbf{Elimination of Calibration Dependency.} 
Our proposed \textbf{AO-FP-CA} matches the performance of the offline-calibrated \textbf{FP-CA} exactly on clean streams (95.23\%) and on contrast-shifted streams (94.92\%--94.96\%). This indicates that pre-computing parameter sensitivities offline on clean, private datasets is completely unavailable. Parameter-level sensitivities can be learned online from unlabeled data on-the-fly, preserving privacy and removing data-sharing requirements.

\textbf{Adaptivity under Corruptions.}
In the absence of prototype routing (LFWA vs. AO-LFWA), our proposed \textbf{AO-LFWA} significantly outperforms offline LFWA in multiple settings. For instance, on alternating clean streams, AO-LFWA achieves \textbf{45.83\%} vs. \textbf{38.56\%} for LFWA. On contrast shift, AO-LFWA achieves \textbf{43.85\%} vs. \textbf{38.25\%}. On sequential Gaussian noise, AO-LFWA reaches \textbf{35.19\%} vs. \textbf{33.73\%}. This shows that pre-computed clean sensitivities are suboptimal under severe test-time noise, whereas estimating sensitivities online captures the actual, noisy loss landscape and provides better-conditioned updates.

\textbf{Trajectory Visualization.}
As shown in \cref{fig:trajectories}, the trajectory of the merging coefficients $\Lambda$ under our active online Fisher estimation tracks the true task boundaries flawlessly. It achieves identical stabilization to the offline-calibrated baseline, demonstrating the technical soundness of online sensitivity updates.

\subsection{Comparison with Weight-Space TTA Paradigms}
To contextualize our contribution, we compare TTMM with classical test-time adaptation (TTA) methods such as TENT \cite{tent}. Classical TTA updates the model parameters $\theta$ directly to minimize entropy. While highly flexible, this weight-space optimization has critical limitations:
\begin{enumerate}
    \item \textbf{High Dimensionality}: Updating millions of parameters on a continuous stream of unlabeled, non-stationary test data is prone to overfitting and representation collapse.
    \item \textbf{Catastrophic Forgetting}: Standard TTA updates the parameters sequentially, which permanently overrides prior expert knowledge. If the stream returns to a previous task, performance collapses.
    \item \textbf{Inability to Merge Experts}: Classical TTA works on a single model and cannot dynamically combine complementary experts (e.g. MNIST and FashionMNIST experts) to solve joint tasks on-the-fly.
\end{enumerate}
By contrast, our proposed \textbf{AO-TTMM} operates in the low-dimensional coefficient space. This regularizes updates, eliminates catastrophic forgetting, and enables collaborative multi-task capabilities, all while remaining completely calibration-free and privacy-preserving.

\subsection{Sensitivity Analysis \& Ablation Studies}
To further analyze the behavior and robustness of our proposed Active Online Fisher Estimation, we conduct systematic hyperparameter sweeps on the key formulation coefficients of AO-FP-CA over the Alternating Stream.

\textbf{Sensitivity to Online Update Rate $\gamma_{\text{fish}}$.}
The parameter $\gamma_{\text{fish}} \in (0,1)$ dictates the rate of information updates in the EMA of squared gradients. Small values of $\gamma_{\text{fish}}$ introduce extreme temporal smoothing, which filters out high-frequency batch noise but delays adaptation to true shifts. Large values update sensitivities aggressively, making estimation susceptible to local batch variances. \cref{tab:gamma_sweep} shows accuracy under varying $\gamma_{\text{fish}}$ values (holding $\alpha = 1.0$).

\begin{table}[t]
\caption{Sensitivity of AO-FP-CA (\%) to the online Fisher update rate $\gamma_{\text{fish}}$ (with $\alpha = 1.0$) on the Alternating Stream.}
\label{tab:gamma_sweep}
\vskip 0.08in
\centering
\begin{small}
\begin{sc}
\begin{tabular}{lccc}
\toprule
$\gamma_{\text{fish}}$ & Clean & Gaussian Noise & Contrast Shift \\
\midrule
0.01 & 95.23\% & 37.96\% & 94.96\% \\
0.05 & 95.23\% & 38.54\% & 94.94\% \\
0.10 & 95.23\% & 38.52\% & 94.96\% \\
0.30 & 95.23\% & 39.48\% & 94.96\% \\
0.50 & 95.23\% & 37.54\% & 94.96\% \\
\bottomrule
\end{tabular}
\end{sc}
\end{small}
\vskip -0.1in
\end{table}

The results show that clean and contrast shift accuracies are exceptionally stable across a wide range of update rates, achieving 95.23\% and 94.96\% respectively. Under severe Gaussian noise, the accuracies are relatively robust, peaking at 39.48\% for $\gamma_{\text{fish}} = 0.30$, while smaller values (e.g., 0.01) maintain 37.96\%, showing that our online Fisher estimation is highly stable across different information update rates.

\textbf{Impact of Preconditioning Sensitivity Exponent $\alpha$.}
The exponent $\alpha$ dictates the intensity of preconditioning. When $\alpha = 0.0$, the preconditioning is disabled, reducing the algorithm to uniform layer-wise learning rates (equivalent to standard AdaMerging with prototypes). Large values of $\alpha$ aggressively slow down sensitive layers. \cref{tab:alpha_sweep} compiles results across different values of $\alpha$ (holding $\gamma_{\text{fish}} = 0.1$).

\begin{table}[t]
\caption{Impact of the preconditioning exponent $\alpha$ (with $\gamma_{\text{fish}} = 0.1$) on AO-FP-CA accuracy (\%) on the Alternating Stream.}
\label{tab:alpha_sweep}
\vskip 0.08in
\centering
\begin{small}
\begin{sc}
\begin{tabular}{lccc}
\toprule
$\alpha$ & Clean & Gaussian Noise & Contrast Shift \\
\midrule
0.0 & 95.23\% & 38.15\% & 94.96\% \\
0.5 & 95.23\% & 40.73\% & 94.94\% \\
1.0 & 95.23\% & 37.29\% & 94.96\% \\
1.5 & 95.23\% & 39.98\% & 94.96\% \\
\bottomrule
\end{tabular}
\end{sc}
\end{small}
\vskip -0.1in
\end{table}

The ablation demonstrates that introducing preconditioning ($\alpha > 0$) leads to a significant performance improvement under Gaussian noise, rising from \textbf{38.15\%} (at $\alpha = 0.0$, i.e., no preconditioning) to a peak of \textbf{40.73\%} (at $\alpha = 0.5$) and \textbf{39.98\%} (at $\alpha = 1.5$). This confirms our core hypothesis: preconditioning the learning rates of sensitive layers is critical under noise, preventing representation collapse and stabilizing online updates.

\section{Practical Considerations for Edge Deployment}
\label{submission-practical-considerations}

Deploying deep learning models on resource-constrained edge devices (e.g., Internet of Things, embedded systems, automotive AI, and mobile devices) imposes strict constraints on both communication bandwidth and memory consumption. In these decentralized environments, communicating full parameter updates or sharing clean calibration datasets is frequently impossible due to communication bottlenecks and strict data-privacy regulations.

Our proposed \textbf{AO-TTMM} is highly optimized for edge deployment. Because it optimizes the low-dimensional merging coefficients $\Lambda$ rather than the full model weights $\theta$, its communication footprint is extremely small (just 63 floats, or 252 bytes, for a ResNet-18 model). This represents an astronomical reduction in bandwidth compared to standard decentralized learning or weight-adaptation methods, which require transmitting millions of weights.

Furthermore, by estimating parameter-level sensitivities on-the-fly directly from incoming unlabeled data streams, AO-TTMM completely avoids the need to store, transmit, or share clean calibration datasets. This provides a robust guarantee of data privacy, as sensitive local user data never leaves the edge device, and no global calibration dataset is required.

\section{Broader Impact and Limitations}
\label{submission-limitations}

While AO-TTMM achieves a complete removal of offline calibration data requirements, it introduces a slight computational overhead during test-time adaptation. Specifically, our method requires computing gradients of the prediction entropy with respect to parameter tensors to update online sensitivities. On edge devices with strict power and speed constraints, this backpropagation step might be expensive. 

To mitigate this, future work can explore: (1) selective layer estimation, where sensitivities are only computed for a subset of highly sensitive layers (e.g., final conv blocks); or (2) periodic estimation, where online Fisher values are computed once every $N$ batches rather than at every step, leveraging the high temporal smoothing stability demonstrated in our ablation studies.

\section{Conclusion}
\label{submission-conclusion}

We presented \textbf{AO-TTMM}, a novel calibration-free approach to test-time model merging that dynamically estimates parameter-level sensitivities from unlabeled test-time data streams. Our experiments show that our online estimation matches or exceeds the performance of offline calibration baselines across various non-stationary streams under severe corruptions. This work represents a significant step toward private, data-free, and adaptive edge intelligence.

\nocite{*}
\bibliography{paper}
\bibliographystyle{icml2026}

\clearpage
\appendix
\section{Theoretical Proofs}
\label{appendix-proofs}

\subsection{Proof of Theorem 3.1}
\label{appendix-proof-thm1}
By unrolling the recurrence relation $F_t^{(l)} = (1 - \gamma) F_{t-1}^{(l)} + \gamma g_t^{(l)}$, we obtain:
\begin{equation}
    F_t^{(l)} = (1 - \gamma)^t F_0^{(l)} + \gamma \sum_{j=1}^t (1 - \gamma)^{t-j} g_j^{(l)}
\end{equation}
where $g_j^{(l)} = \text{mean}((g^{(l)})^2)$ at step $j$. Taking the expectation:
\begin{equation}
    \mathbb{E}[F_t^{(l)}] = (1 - \gamma)^t F_0^{(l)} + \gamma \left( \sum_{j=1}^t (1 - \gamma)^{t-j} \right) F_*^{(l)}
\end{equation}
Evaluating the geometric series $\sum_{j=1}^t (1 - \gamma)^{t-j} = \frac{1 - (1 - \gamma)^t}{\gamma}$ yields:
\begin{equation}
    \mathbb{E}[F_t^{(l)}] = (1 - \gamma)^t F_0^{(l)} + (1 - (1 - \gamma)^t) F_*^{(l)}
\end{equation}
Since $|1 - \gamma| < 1$, taking $t \to \infty$ results in $\lim_{t \to \infty} \mathbb{E}[F_t^{(l)}] = F_*^{(l)}$, proving asymptotic unbiasedness.

For the variance, since $F_0^{(l)}$ is constant and the batches are i.i.d.:
\begin{align}
    \text{Var}(F_t^{(l)}) &= \text{Var}\left( \gamma \sum_{j=1}^t (1 - \gamma)^{t-j} g_j^{(l)} \right) \\
    &= \gamma^2 \sum_{j=1}^t (1 - \gamma)^{2(t-j)} \text{Var}(g_j^{(l)})
\end{align}
Using the variance bound $\text{Var}(g_j^{(l)}) \leq \sigma^2$:
\begin{align}
    \text{Var}(F_t^{(l)}) &\leq \gamma^2 \sigma^2 \sum_{k=0}^{t-1} (1 - \gamma)^{2k} \\
    &= \gamma^2 \sigma^2 \frac{1 - (1 - \gamma)^{2t}}{1 - (1 - \gamma)^2}
\end{align}
Simplifying the denominator $1 - (1 - \gamma)^2 = \gamma(2 - \gamma)$ gives:
\begin{equation}
    \text{Var}(F_t^{(l)}) \leq \frac{\gamma (1 - (1 - \gamma)^{2t})}{2 - \gamma} \sigma^2
\end{equation}
Taking the limit as $t \to \infty$ yields $\lim_{t \to \infty} \text{Var}(F_t^{(l)}) \leq \frac{\gamma}{2 - \gamma} \sigma^2$.

\subsection{Proof of Theorem 3.2}
\label{appendix-proof-thm2}
Taking the expectation of the batch-wise estimator $g_B^2$:
\begin{align}
    \mathbb{E}[g_B^2] &= \mathbb{E}\left[ \left( \frac{1}{B} \sum_{i=1}^B g_i \right)^2 \right] \\
    &= \frac{1}{B^2} \sum_{i=1}^B \mathbb{E}[g_i^2] + \frac{1}{B^2} \sum_{i \neq j} \mathbb{E}[g_i g_j] \\
    &= \frac{1}{B^2} \left( B (\mu^2 + \sigma^2) + B(B-1) \mu^2 \right) \\
    &= \mu^2 + \frac{\sigma^2}{B}
\end{align}
Subtracting the true single-sample expected squared gradient $F_* = \mu^2 + \sigma^2$ yields the bias:
\begin{equation}
    \text{Bias}(g_B^2) = \mu^2 + \frac{\sigma^2}{B} - (\mu^2 + \sigma^2) = - \frac{B-1}{B} \sigma^2
\end{equation}
For the variance, the sample mean $g_B = \frac{1}{B}\sum_{i=1}^B g_i \sim \mathcal{N}(\mu, \frac{\sigma^2}{B})$. For a normal random variable $X \sim \mathcal{N}(m, s^2)$, the variance of its square is $\text{Var}(X^2) = 2 s^4 + 4 m^2 s^2$. Substituting $m = \mu$ and $s^2 = \frac{\sigma^2}{B}$ yields:
\begin{equation}
    \text{Var}(g_B^2) = 2 \left( \frac{\sigma^2}{B} \right)^2 + 4 \mu^2 \left( \frac{\sigma^2}{B} \right) = \frac{2 \sigma^4}{B^2} + \frac{4 \mu^2 \sigma^2}{B}
\end{equation}

\end{document}
